\begin{theorem}[Canonical representative of an injective MPU]
    \label{thm:mpu_injective_cfii_representative}
    \lean{MPOTensor.IsMPU.exists_cfii_representative_of_isInjective}
    \leanok
    \uses{def:mpu, def:injective, def:mpu_canonical_form_ii}
    Let $d>0$ and let $\mathcal U$ be an injective MPU tensor of bond
    dimension $D$. Combining the canonical replacement of
    \cite[lines~319--356]{Cirac2017MPU} with the equality case of
    \cite[Proposition~20]{Molnar2018SemiInjective}, there are a
    canonical-form-II tensor $\mathcal V$ of the same bond dimension and an
    invertible matrix $X$ such that, for every $N>0$,
    \begin{align}
        \mathcal V^{ij} & =X\mathcal U^{ij}X^{-1}, \\
        V^{(N)}         & =U^{(N)}.
    \end{align}
\end{theorem}

\begin{proof}\leanok
    \uses{thm:mpu_reduced_cfii_representative,
        thm:mps_equality_reduction_exists, thm:mps_rectangular_reduction_identities}
    A reduced canonical representative $\mathcal V$ has bond dimension
    $E\leq D$, and equality of the positive-length periodic operators
    identifies the doubled-index matrix product vectors of $\mathcal U$ and
    $\mathcal V$. Since $\mathcal U$ is injective, there are
    $P\in\C^{D\times E}$ and $Q\in\C^{E\times D}$ with $PQ=I_D$ and
    $P\mathcal V^{ij}Q=\mathcal U^{ij}$. This left-invertibility implies
    $D\leq E$ and hence $D=E$, so for square matrices $PQ=I_D$ also gives
    $QP=I_D$. Multiplying the letter identity on the left by $Q$ and on the
    right by $P$ yields $\mathcal V^{ij}=Q\mathcal U^{ij}P$, as required
    with $X=Q$.
\end{proof}

\begin{theorem}[Source ranks and index of a simple injective MPU]
    \label{thm:mpu_simple_injective_source_index}
    \lean{MPOTensor.IsMPU.rightRank_mul_leftRank_of_isInjective_of_isMPUSimple,
        MPOTensor.IsMPU.index_eq_logb_of_isInjective_of_isMPUSimple}
    \leanok
    \uses{def:mpu, def:injective, def:mpu_simple,
        def:mpu_reduced_representative_index, def:mpu_right_rank, def:mpu_left_rank}
    Let $d>0$ and let $\mathcal U$ be a simple injective MPU tensor as in
    \cite[lines~1403--1407]{FrancoRubio2025MPUGauging}. Following
    \cite[lines~562--573 and~681--704]{Cirac2017MPU}, its source-cut ranks
    $r$ and $\ell$ satisfy
    \begin{align}
        r\ell & =d^2,                                    \\
        \operatorname{ind}(\mathcal U)
              & =\frac12\bigl(\log_2 r-\log_2\ell\bigr).
    \end{align}
\end{theorem}

\begin{proof}\leanok
    \uses{thm:mpu_injective_cfii_representative,
        thm:mpu_simple_rectangular_inclusion, thm:mpu_virtual_sandwich_source_ranks,
        thm:mpu_reduced_representative_index_invariance,
        thm:mpu_bounded_simple_blocking,
        thm:mpu_source_ranks_physical_reindexing,
        ThmFund1, prop:mpu_admissible_index_well_defined}
    Passing to the same-dimensional canonical representative
    $\mathcal V^{ij}=X\mathcal U^{ij}X^{-1}$ preserves both simplicity and
    the two source ranks. The simple canonical tensor $\mathcal V$ therefore
    has rank product $r\ell=d^2$. Since its one-site blocking only relabels
    the physical alphabet, its index is
    $\frac12(\log_2 r-\log_2\ell)$, and equality of the periodic operators
    identifies this value with $\operatorname{ind}(\mathcal U)$.
\end{proof}
